\section{Metrics}
\label{sec:metrics}

Two things are reported for every policy: what the model achieved, and what it cost to get
there. The second is the study's subject. A policy that reaches a good model slowly has spent
rounds its facilities could not spare. Which quality metric the first of those uses is not fixed
here; that choice belongs to the review of related literature, which does not yet exist
(Section~\ref{sec:open}).

\subsection{Notation}

\begin{spectab}{|L{0.16\textwidth}|L{0.74\textwidth}|}
\hline
\rowcolor{headcell}\textbf{Symbol} & \textbf{Meaning} \\ \hline
\endfirsthead
\hline
\rowcolor{headcell}\textbf{Symbol} & \textbf{Meaning} \\ \hline
\endhead
$i$ & facility index \\ \hline
\rowcolor{rowshade}$r$ & round index, counted from 1 \\ \hline
$c_i(r)$ & number of rounds facility $i$ contributed to, after round $r$ \\ \hline
\rowcolor{rowshade}$P_r$ & facilities that actually trained in round $r$ \\ \hline
$\ell_i$ & time facility $i$ needs to complete one round, in seconds \\ \hline
\rowcolor{rowshade}$d_r$ & duration of round $r$, in seconds \\ \hline
$t_r$ & simulated elapsed time after $r$ rounds \\ \hline
\rowcolor{rowshade}$Q_r$ & the reported quality metric after $r$ rounds (Section~\ref{sec:open}) \\ \hline
$\tau$ & target value of $Q_r$ for a usable model \\ \hline
\rowcolor{rowshade}$r^{*}$ & the first round at which $Q_r$ meets $\tau$ \\ \hline
$T^{*}$ & the elapsed time at that round, $t_{r^{*}}$ \\ \hline
\rowcolor{rowshade}$R_p$ & the first round at which every facility has met the floor $p$ \\ \hline
$p$ & participation floor, as a share of rounds \\ \hline
\rowcolor{rowshade}$h_i$ & rows held out at facility $i$ (Section~\ref{sec:quality}) \\ \hline
\end{spectab}

\noindent
Throughout, a star marks the round at which a condition first holds: notation, never a
multiplication. Both minima below range over the
$R$ rounds of a single run (Section~\ref{sec:parameters}), so a condition met in no round of the
run leaves its quantity undefined. The run's own parameters --- the round budget $R$, the
held-out fraction $f$ and the record volume $n_i$ --- are listed in
Section~\ref{sec:parameters}; within a sum, $j$ indexes a round. In the listings,
\texttt{c[i][r]}, \texttt{d[r]} and \texttt{t[r]} are $c_i(r)$, $d_r$ and $t_r$. Three
conventions apply: time is
simulated rather than wall-clock, because the study prices time and not compute; every figure
carries its seed, configuration and data version; and a result that cannot be computed is printed
as undefined rather than coerced to zero.

\subsection{Model quality}
\label{sec:quality}

The study classifies each facility period as a surge or not, from facility records, and two ways
of scoring such a model are defensible. The period can be treated as a quantity and scored with a
scaled error measure, or the surge label can be treated as a class and scored with balanced
accuracy, the mean per-class recall that a rare class makes necessary. They answer different
questions. An error measure says how far a count is missed by; balanced accuracy says how reliably
a surge is caught.

The study chooses balanced accuracy on surge periods (OD-7, resolved 2026-11-20). Two reasons. A
low-volume facility's demand series is zero-heavy, and a percentage error is undefined or explosive
on zeros, whereas balanced accuracy never divides by a count that may be zero. And the comparison
this study reports is per facility: an error measure is not scale-free, so ranking a station
holding a few hundred records against one holding a few thousand would need a scaling argument
before any ranking could be stated. Balanced accuracy also has a floor a reader can reason with ---
a trivial always-negative classifier scores $0.5$ --- which is what makes a target interpretable.

\noindent
Balanced accuracy is the mean of the per-class recall rates, so that a rare class carries the same
weight as a common one:

\begin{equation}
\mathrm{BA} \;=\; \frac{1}{K}\sum_{k=1}^{K}
  \frac{\text{rows of class } k \text{ classified correctly}}{\text{all rows of class } k},
\label{eq:ba}
\end{equation}

\noindent
with $K$ classes and, here, $K = 2$. On a pooled held-out set of $100$ surge rows and $400$
non-surge rows, a model that catches $60$ surges and classifies $380$ of the non-surges correctly
scores

\[
\mathrm{BA} \;=\; \frac{1}{2}\left(\frac{60}{100} + \frac{380}{400}\right)
\;=\; \frac{0.60 + 0.95}{2} \;=\; 0.775 .
\]

\noindent
Plain accuracy cannot stand in. Those same predictions give $(60 + 380)/500 = 0.88$ plain accuracy,
but a model that always answers ``no surge'' gives $400/500 = 0.80$ plain accuracy while catching no
surge at all, and only $0.50$ balanced accuracy. Averaging the two recall rates is what keeps a rare
event from being averaged away.

What the document fixes is the shape of the
requirement: one scalar $Q_r$, comparable across policies, computed on the pooled held-out rows
of every facility, split chronologically because neighbouring clinic-days are predictive, and
carrying its own direction --- for balanced accuracy, a rate, so higher is better.

\noindent
The evaluation itself is fixed:

{\small
\begin{verbatim}
for each facility i:
    sort the rows of i by period, oldest first
    h <- ceil(f * n[i])            # f from Section 7
    held-out[i] <- last h rows of i    # none if h = 0
    training[i] <- the rows before them
pooled <- union of held-out rows over all i
valid <- pooled non-empty and all training[i] non-empty
for r = 1 to R:
    if valid:
        Q[r] <- score(model after round r, pooled)
    else:
        Q[r] <- undefined          # nothing comparable to score
output Q[1..R]
\end{verbatim}
}

\noindent
Each facility holds out the last $h_i$ rows of its own record, where

\begin{equation}
h_i \;=\; \lceil f\, n_i \rceil
\label{eq:holdout}
\end{equation}

\noindent
and the pooled evaluation set is the sum of those counts:

\begin{equation}
\text{pooled rows} \;=\; \sum_i h_i
\label{eq:pooled}
\end{equation}

\noindent
In arithmetic, at $f = 0.2$ and $n_i = 350$:

\[
h_i \;=\; \lceil 0.2 \times 350 \rceil \;=\; 70\ \text{held out},
\qquad
350 - 70 \;=\; 280\ \text{trained on},
\]

\noindent
and three facilities of that size give a pooled evaluation set of

\[
\sum_i h_i \;=\; 70 + 70 + 70 \;=\; 210\ \text{rows}.
\]

\noindent
A configuration
that leaves a facility with no training rows, or leaves the pooled set empty, has no comparable
quality: $Q_r$ is then undefined for every round, and the run is reported as invalid rather than
scored.

\noindent\textbf{In plain terms.} Quality is scored only on records the model never saw while
training, so it cannot look good by memorising. The two options answer different questions: how
far short a count may be missed by, or how reliably a surge period is flagged. A stock decision,
or a warning.

\subsection{Time to a usable model}

The clock starts when training starts and stops at the round where quality is good enough. Whether
a value clears the target depends on which way the metric runs:

\[
Q_r \text{ meets } \tau \quad\Longleftrightarrow\quad
\begin{cases}
Q_r \le \tau, & \text{when } Q \text{ is an error},\\[2pt]
Q_r \ge \tau, & \text{when } Q \text{ is a rate}.
\end{cases}
\]

\begin{equation}
r^{*} = \min\{\, r : Q_r \text{ meets } \tau \,\},
\qquad T^{*} = t_{r^{*}}
\label{eq:tstar}
\end{equation}

\noindent
If no round reaches $\tau$, $T^{*}$ is \emph{undefined}, not zero. A policy that never produces
a usable model has failed in a way a zero would hide.

\noindent
$\tau$ is not chosen from outside the study. It is the balanced accuracy of the local-only baseline
--- a model trained on one facility's own records --- computed once per experimental condition from
the frozen baselines before any policy comparison runs, and declared with that condition
(Section~\ref{sec:open}, OD-7). $T^{*}$ then answers something a facility can act on: how long
until federation catches up to what it could already do alone. A condition whose local-only
baseline falls at or below $0.5$ has no usable $\tau$, since $0.5$ is what an always-negative
classifier scores; such a condition is reported as unevaluable rather than compared.

\noindent
The same definition as a procedure, for readers who find the set notation harder than the loop:

{\small
\begin{verbatim}
time <- 0
r* <- undefined
T* <- undefined
for r = 1 to R:                # the run is R rounds (Section 7)
    d[r] <- round duration      # from 6.5; 0 if nobody answered
    time <- time + d[r]        # after this line, time = t_r
    if Q[r] is defined and Q[r] meets tau:    # <= error, >= rate
        r* <- r
        T* <- time
        break
output (r*, T*)                # undefined pair if none met tau
\end{verbatim}
}

\noindent
The same quantity in arithmetic: $T^{*}$ is the sum of the round durations up to and including
the round that reaches the target.

\begin{equation}
T^{*} \;=\; d_1 + d_2 + \cdots + d_{r^{*}} \;=\; \sum_{j=1}^{r^{*}} d_j
\label{eq:tstar-sum}
\end{equation}

\noindent
If the target is first reached at round four, and those four rounds took two, five, three and
four hours, then

\[
T^{*} \;=\; 2 + 5 + 3 + 4 \;=\; 14\ \text{hours}.
\]

\noindent
The durations are quoted in hours here for readability; the model itself counts seconds. A round
that stalls contributes its zero like any other term in the sum.

\noindent\textbf{In plain terms.} Training runs in rounds, and each round costs time: staff, power,
a working connection. Quality does not rise evenly, so what matters is when the model becomes good
enough to trust. $r^{*}$ is that round, and $T^{*}$ is the clock when it arrives.

\subsection{Participation floor and its cost}

\begin{equation}
R_p = \min\left\{\, r : \min_i \frac{c_i(r)}{r} \;\ge\; p \,\right\},
\qquad \text{cost} = t_{R_p}
\label{eq:rfloor}
\end{equation}

\noindent
$R_p$ is the first round at which every facility has met the floor \emph{at the same time}.
Because the share is cumulative, a facility that reaches it once can fall below it again by
missing later rounds, so $R_p$ is read alongside each facility's trajectory $c_i(r)/r$; the
single number would otherwise claim a sustained inclusion the data has not shown. A facility that
never reaches the floor leaves $R_p$ undefined: that is the finding, not a missing value.

\noindent
The search for that round, as a procedure:

{\small
\begin{verbatim}
Rp <- undefined
cost <- undefined
for r = 1 to R:
    low <- min over all i of c[i][r]/r    # every facility
    if low >= p:
        Rp <- r
        cost <- t[r]
        break
output (Rp, cost)
\end{verbatim}
}

\noindent
Its cost is the same elapsed-time sum as any other round:

\begin{equation}
t_{R_p} \;=\; d_1 + d_2 + \cdots + d_{R_p} \;=\; \sum_{j=1}^{R_p} d_j
\label{eq:rfloor-cost}
\end{equation}

\noindent
The round itself has an arithmetic form too. Because $c_i(r)$ is a whole number,

\[
\frac{c_i(r)}{r} \;\ge\; p \quad\Longleftrightarrow\quad c_i(r) \;\ge\; \lceil p\,r \rceil,
\]

\noindent
so $R_p$ is the smallest $r$ at which the right-hand inequality holds for every facility. Take
$p = 0.75$ and compare six rounds; each facility must then have joined at least

\[
\lceil 0.75 \times 6 \rceil \;=\; \lceil 4.5 \rceil \;=\; 5
\]

\noindent
of the six. If those six rounds took three, four, two, five, three and two hours, the cost is

\[
t_{R_p} \;=\; 3 + 4 + 2 + 5 + 3 + 2 \;=\; 19\ \text{hours}.
\]

\noindent\textbf{In plain terms.} The floor exists to catch a health center being quietly left out:
it measures whether every facility took part in at least a stated share of rounds, so that the
model in use is built with each of them, not around them. $R_p$ is the first round where all have
reached that share at once, and the cost is how long it took. Reaching it once is not holding it.

\subsection{Round duration and elapsed time}

\begin{equation}
d_r = \max_{i \in P_r} \ell_i,
\qquad t_r = \sum_{j=1}^{r} d_j
\label{eq:duration}
\end{equation}

\noindent
$\ell_i$ is the whole per-round cost to facility $i$: local training, then transfer in both
directions. A synchronous round ends when the last update arrives, so $d_r$ is the slowest
participant's time and not the average, and one slow facility sets the pace for every other
participant in that round. Server-side work (selection and aggregation) is assumed to take no
time: what is priced is the facilities' time, not the aggregator's.

The maximum in equation~\eqref{eq:duration} is taken as zero when no participant answers, so a
stalled round has $d_r = 0$ and no progress claimed. That zero is a modelling assumption, not a
measurement: a deployment with timeouts would charge the timeout instead, and stating the
assumption keeps the difference visible.

\noindent
Round by round, those two equations are one loop:

{\small
\begin{verbatim}
t[0] <- 0
for r = 1 to R:
    if P[r] is empty:
        d[r] <- 0                  # stalled round
    else:
        d[r] <- max ell[i], i in P[r]
    t[r] <- t[r-1] + d[r]
output t[1..R], d[1..R]
\end{verbatim}
}

\noindent
\[
t_r \;=\; d_1 + d_2 + \cdots + d_r
\]

\noindent
Three rounds lasting two, zero and three hours, with the middle one stalled, give

\[
t_3 \;=\; 2 + 0 + 3 \;=\; 5\ \text{hours}.
\]

\noindent\textbf{In plain terms.} A round is not finished when the fastest center answers. It is
finished when the slowest one does, because the server waits for every update before it can
combine them. One facility having a bad week therefore sets the pace for everyone that week, and
the total time is those rounds added up.
